\usepackage[T1]{fontenc}        % codificación de fuente para fuentes vectoriales
\usepackage{newtxtext,newtxmath} % tipografía Times compatible con las opciones APA 7
\usepackage[spanish]{babel}     % español
\AtBeginDocument{\shorthandoff{<>}} % evita que babel-spanish active "<" ">" (rompe flechas de TikZ)
\usepackage[utf8]{inputenc}     % acentos sin codigo
\usepackage[margin=2.54cm,headsep=0.5cm,footskip=0.8cm]{geometry}
\usepackage{setspace}
\usepackage{titlesec}
\usepackage{amsmath} % Paquetes matemáticos de la American Mathematical Society; newtxmath ya provee sus símbolos
\usepackage{graphicx}		  % para imagenes
\usepackage{float} 		      % para imagenes en lugar solicitado H
\usepackage{listings}	      % para codigo fuente
\usepackage[shortlabels]{enumitem} % para listas personalizadas y glosarios
\usepackage{url}			  % referencias bibliograficas WEB
\usepackage{xurl}             % corte inteligente de URLs largas
\usepackage{microtype}        % optimización tipográfica del espaciado
\usepackage{ragged2e}         % alineacion de texto
\usepackage{caption}          % formato APA 7 de títulos de tablas y figuras
\usepackage{subcaption}	      % para imagenes multiples
\usepackage{indentfirst}      % sangría también en el primer párrafo tras un título
\usepackage{fancyhdr}         % folio en la esquina superior derecha
\usepackage{booktabs}         % Líneas elegantes en tablas
\usepackage{threeparttable}   % Notas al pie en tablas
\usepackage{tabularx}         % Tablas con ancho controlado (\textwidth)
\usepackage{array}            % Modificadores de columna (>{...})
\newcolumntype{P}[1]{>{\RaggedRight\arraybackslash}p{#1}} % texto alineado a la izquierda
\newcolumntype{L}{>{\RaggedRight\arraybackslash}X}       % columna flexible alineada a la izquierda
\usepackage{icomma}           % Manejo de coma decimal en modo matemático
\usepackage[table]{xcolor}    % paleta de color y color en tablas (\rowcolor)
\usepackage{tikz}             % diagramas vectoriales
\usetikzlibrary{shapes.geometric, arrows.meta, positioning, fit, backgrounds, calc}
\usepackage{pgfplots}         % gráficos de datos vectoriales
\pgfplotsset{compat=1.18}
\usepackage{pgfgantt}         % diagrama de Gantt vectorial
\usepackage{chngcntr}
\usepackage{csquotes}
\usepackage[backend=biber,style=apa,sorting=nyt]{biblatex}
\DeclareLanguageMapping{spanish}{spanish-apa}
\addbibresource{partes/Bibliografia.bib}
\setlength{\bibhang}{1.27cm} % sangría francesa APA 7
\setlength{\bibitemsep}{0pt}
\usepackage{hyperref} 		  % para generar marcadores en el PDF ponerlo antes del \begin{document}
\usepackage{etoolbox}
\providecommand{\tightlist}{\setlength{\itemsep}{0pt}\setlength{\parskip}{0pt}}

% ---- Paleta de color unificada para los diagramas del plan de tesis ----
% Paleta del título de tesis, compartida con la presentación de formulación.
\definecolor{bpBlue}{RGB}{43,91,132}
\definecolor{bpGold}{RGB}{200,157,66}
\definecolor{bpRust}{RGB}{179,92,62}

\definecolor{tcAzul}{HTML}{2E5C8A}    % elementos centrales / cliente / orquestación
\definecolor{tcTeal}{HTML}{15807A}    % procesos / recuperación / consulta
\definecolor{tcNaranja}{HTML}{C9702B} % modelo / inferencia / configuraciones E1-E3
\definecolor{tcVerde}{HTML}{3F8F52}   % datos / corpus / almacenamiento / fines
\definecolor{tcMorado}{HTML}{6A4C93}  % evaluación / análisis / síntesis
\definecolor{tcRojo}{HTML}{B33A3A}    % efectos negativos / riesgos
\definecolor{tcGris}{HTML}{5B6470}    % bordes y flechas neutras

\tikzset{
  tcaja/.style={draw=tcGris, rounded corners=2.5pt, align=center, font=\scriptsize,
    line width=0.7pt, inner sep=4pt},
  tcentral/.style={tcaja, fill=tcAzul!14, draw=tcAzul, line width=1.1pt, font=\scriptsize\bfseries, text=tcAzul!55!black},
  tazul/.style={tcaja, fill=tcAzul!12, draw=tcAzul!75!black},
  tteal/.style={tcaja, fill=tcTeal!14, draw=tcTeal!80!black},
  tnaranja/.style={tcaja, fill=tcNaranja!16, draw=tcNaranja!85!black},
  tverde/.style={tcaja, fill=tcVerde!14, draw=tcVerde!75!black},
  tmorado/.style={tcaja, fill=tcMorado!14, draw=tcMorado!80!black},
  trojo/.style={tcaja, fill=tcRojo!12, draw=tcRojo!80!black},
  techa/.style={-{Latex[length=2.2mm]}, thick, draw=tcGris},
  tlbl/.style={font=\scriptsize\itshape, text=tcGris}
}

% Tablas y figuras estilo APA 7: etiqueta en negrita y título en cursiva,
% cada uno en una línea propia y alineados a la izquierda.
\captionsetup[figure]{labelfont={bf}, textfont={it} ,labelformat={default}, labelsep=newline, name={Figura }, justification=raggedright,singlelinecheck=false} %-----formato del caption figure
\captionsetup[table]{labelfont={bf}, textfont={it} ,labelformat={default}, labelsep=newline, name={Tabla },justification=raggedright,singlelinecheck=false} %-----formato del caption table

\hypersetup{hidelinks,
		colorlinks=true,
		allcolors=black,
		pdfstartview=Fit,
		breaklinks=true,
		hypertexnames=false
	}	%para quitar marco rojo y evitar identificadores de pagina duplicados

% ---- Configuración de flotante para Algoritmos ----
\newfloat{algoritmo}{H}{loa}
\floatname{algoritmo}{Algoritmo}
\captionsetup[algoritmo]{labelfont={bf}, textfont={it}, labelformat={default}, labelsep=newline, name={Algoritmo }, justification=raggedright, singlelinecheck=false}

% Folio APA 7 en el encabezado derecho, también en páginas de capítulo.
\setlength{\headheight}{15pt}
\pagestyle{fancy}
\fancyhf{}
\fancyhead[R]{\thepage}
\renewcommand{\headrulewidth}{0pt}
\fancypagestyle{plain}{%
  \fancyhf{}%
  \fancyhead[R]{\thepage}%
  \renewcommand{\headrulewidth}{0pt}%
}

% Capítulos romanos y apartados numerados según el esquema institucional.
\renewcommand{\thechapter}{\Roman{chapter}}
\renewcommand{\thesection}{\arabic{chapter}.\arabic{section}}
\renewcommand{\thesubsection}{\thesection.\arabic{subsection}}
\renewcommand{\thesubsubsection}{\thesubsection.\arabic{subsubsection}}
\setcounter{secnumdepth}{5}
\setcounter{tocdepth}{5}

% Jerarquía APA 7 a 12 pt. Se mantiene la identificación institucional
% “Capítulo I” y la numeración del plan; los cinco niveles siguen el formato APA.
\titleformat{\chapter}[block]
  {\normalfont\normalsize\bfseries\centering}
  {\MakeUppercase{\chaptername}~\thechapter}{0.5em}{\MakeUppercase}
\titlespacing*{\chapter}{0pt}{0pt}{\baselineskip}
\titleformat{\section}[block]
  {\normalfont\normalsize\bfseries\centering}{\thesection}{0.5em}{}
\titlespacing*{\section}{0pt}{\baselineskip}{0pt}
\titleformat{\subsection}[block]
  {\normalfont\normalsize\bfseries}{\thesubsection}{0.5em}{}
\titlespacing*{\subsection}{0pt}{\baselineskip}{0pt}
\titleformat{\subsubsection}[block]
  {\normalfont\normalsize\bfseries\itshape}{\thesubsubsection}{0.5em}{}
\titlespacing*{\subsubsection}{0pt}{\baselineskip}{0pt}
\titleformat{\paragraph}[runin]
  {\normalfont\normalsize\bfseries}{\theparagraph}{0.5em}{}[.]
\titlespacing*{\paragraph}{1.27cm}{\baselineskip}{0.5em}
\titleformat{\subparagraph}[runin]
  {\normalfont\normalsize\bfseries\itshape}{\thesubparagraph}{0.5em}{}[.]
\titlespacing*{\subparagraph}{1.27cm}{\baselineskip}{0.5em}

% Numeración consecutiva de tablas, figuras y ecuaciones, como en APA 7.
\counterwithout{figure}{chapter}
\counterwithout{table}{chapter}
\counterwithout{equation}{chapter}
\renewcommand{\thefigure}{\arabic{figure}}
\renewcommand{\thetable}{\arabic{table}}
\renewcommand{\theequation}{\arabic{equation}}

% Texto a doble espacio, alineado a la izquierda y con sangría APA 7.
% Para justificarlo por exigencia institucional, sustituir \raggedright por \justifying.
\AtBeginDocument{\doublespacing\raggedright\setlength{\parindent}{1.27cm}}
% El interlineado simple se limita al cuerpo de las tablas y los diagramas;
% los títulos y notas conservan el interlineado general del documento.
\AtBeginEnvironment{tabular}{\singlespacing}
\AtBeginEnvironment{tabularx}{\singlespacing}
\AtBeginEnvironment{tikzpicture}{\singlespacing}
\AtBeginEnvironment{tablenotes}{\fontsize{10}{12}\selectfont\singlespacing}
\AtBeginBibliography{\doublespacing}
